
\subsubsection{Benchmark Contamination Detection}

N-gram overlap methods identify exact or near-exact matches between training and test examples. Recent work recommends n=8 as the detection threshold. Membership inference approaches (Min-K\%++, CDD) operate without explicit training corpus access. The Contamination Taxonomy categorizes contamination by type and impact mechanism. These methods detect contamination presence but do not explain how curation decisions affect contamination rates.

\subsubsection{Data Attribution Methods}

Influence functions approximate leave-one-out effects but scale poorly to large models. TRAK enables practical attribution through random projection, achieving substantial speedups via the LoGra algorithm. Attribution has been applied to data valuation and detecting mislabeled examples, but has not been combined with contamination detection to trace which influential examples derive from benchmark overlap.

\subsubsection{Training Data Curation}

DataComp-LM established that model-based filtering significantly affects downstream performance. SlimPajama demonstrated that deduplication strategies affect model capabilities. FineWeb pipelines combine multiple filtering stages for web-scale curation. Perplexity-based filtering---selecting documents that score well against reference language models---has emerged as a core technique. However, curation research has not analyzed contamination effects.
